\documentclass[11pt]{article}
\usepackage[margin=1in]{geometry}
\usepackage{amsmath,amssymb,amsthm,fontspec,fancyvrb}
\usepackage[title]{appendix}
\setmonofont{DejaVu Sans Mono}[Scale=0.8]
\usepackage[colorlinks=true,linkcolor=blue,citecolor=blue,urlcolor=blue]{hyperref}
\hypersetup{pdftitle={A conjecture of Hopkins on golden Rado numbers},pdfauthor={P. Graham}}
\newtheorem{theorem}{Theorem}
\newtheorem{lemma}[theorem]{Lemma}
\theoremstyle{remark}
\newtheorem{remark}[theorem]{Remark}
\newcommand{\Rg}{R_0^{\varphi}}
\title{A conjecture of Hopkins on golden Rado numbers}
\author{P. Graham\\ \small\url{https://orcid.org/0009-0007-8493-1505}}
\date{October 31, 2026}
\begin{document}
\maketitle
\begin{abstract}
A 2-coloring of $\{1,\dots,N\}$ satisfies the golden mean shift if no two
consecutive integers both receive color~1. For $1\le a\le b$, Hopkins
defined the golden Rado number $\Rg(a,b)$ as the least $N$ such that
every such coloring with $1$ colored~$0$ contains a monochromatic
solution of $ax+by=z$, and conjectured that
$\Rg(a,\ell a)=a^2(\ell+1)^2$ for all positive integers $a$ and $\ell$.
Hopkins proved the case $\ell=1$, and determined $\Rg(1,b)$ for every
$b$; that evaluation shows the formula fails when $a=1$ and $\ell\ge2$.
We prove the formula for every $a\ge2$ and $\ell\ge1$, which settles all
the remaining cases of the conjecture. The lower bound is the coloring
that gives color~1 exactly to the multiples of $a$ that are at least
$a(\ell+1)$. The upper bound is a forcing argument: a descent in steps
of $\ell$ shows that a residue $r\le\min(a,\ell)$ has color~1, and two
families of solutions then make color~1 invariant under shifts by $\ell$
on an initial interval, which leads to a contradiction. The proof is
elementary, and the theorem, both bounds included, is checked in the
Lean proof assistant from the definitions.

This work has not yet been reviewed by experts. The proofs, the Lean
formalization and the text were produced by AI models (from OpenAI and
Anthropic), working under the author's direction.
\end{abstract}
\medskip

\noindent 2020 \emph{Mathematics Subject Classification}: 05D10, 11B75.\\
\emph{Keywords}: Rado number, golden mean shift, Schur-type problem,
Lean.

\section{Introduction}

\subsection{Main result}
A \emph{golden coloring} of $[1,N]=\{1,\dots,N\}$ is a map
$\Delta:[1,N]\to\{0,1\}$ with $\Delta(1)=0$ and no $i$ with
$\Delta(i)=\Delta(i+1)=1$. For positive integers $a,b$, a solution
$(x,y,z)$ of $ax+by=z$ in $[1,N]$ is \emph{monochromatic} if
$\Delta(x)=\Delta(y)=\Delta(z)$; the entries $x$ and $y$ may be equal.
Following Hopkins~\cite[Definition~2]{Hopkins}, for $1\le a\le b$ the
golden Rado number $\Rg(a,b)$ is the least positive integer $N$ such
that every golden coloring of $[1,N]$ contains a monochromatic solution
of $ax+by=z$. (Hopkins writes $R_0^S$ for a binary shift $S$ and
$\Rg$ for the golden mean shift; the subscript records $\Delta(1)=0$.)
For example, $\Rg(1,1)=4$~\cite[Section~2]{Hopkins}. Suppose a golden
coloring of $[1,4]$ has no monochromatic solution of $x+y=z$. Then
$\Delta(2)=1$ because of $(1,1,2)$, so $\Delta(3)=0$, and either color
of $4$ completes $(1,3,4)$ or $(2,2,4)$, a contradiction; the coloring
$0,1,0$ of $[1,3]$ avoids all solutions.

\begin{theorem}\label{thm:main}
For all integers $a\ge2$ and $\ell\ge1$,
\[
\Rg(a,\ell a)=a^2(\ell+1)^2 .
\]
\end{theorem}

Hopkins stated the formula as follows~\cite[Conjecture~1]{Hopkins}:
``Given positive integers $a$ and $\ell$, the golden Rado number for the
case $\Delta(1)=0$ is $\Rg(a,\ell a)=(\ell+1)^2a^2$.'' The case $\ell=1$
is Hopkins's Theorem~2, $\Rg(a,a)=4a^2$~\cite[Theorem~2]{Hopkins}.
For $a=1$, Hopkins's Theorem~1 gives $\Rg(1,b)=2b+2$ for odd $b$ and
$\Rg(1,b)=2b$ for even $b$~\cite[Theorem~1]{Hopkins}; this equals
$(b+1)^2$ only for $b=1$, so the conjectured formula holds at $a=1$ only
when $\ell=1$. Theorem~\ref{thm:main} proves it in every other case.

\subsection{Background}
Without the coloring constraint, the 2-color Rado number of $ax+by=z$
for $1\le a\le b$ is $a(a+b)^2+b$, as conjectured by Hopkins and
Schaal~\cite{HopkinsSchaal} and proved by Guo and Sun~\cite{GuoSun} (we
quote the formula as stated in~\cite[equation~(1)]{Hopkins}). Hopkins
introduced shift-constrained Rado numbers in~\cite{Hopkins}, where the
colorings are required to be words of a binary shift from symbolic
dynamics~\cite{LindMarcus}; the golden mean shift forbids two
consecutive ones. Since the color of $1$ can no longer be fixed by
symmetry, Hopkins treats $\Delta(1)=0$ and $\Delta(1)=1$ separately and
writes that ``a full understanding of $R_0^\varphi(a,b)$,
$R_1^\varphi(a,b)$ is desired''~\cite[Section~5]{Hopkins}. Hopkins
supported Conjecture~1 with computations and acknowledges earlier
exploration of these topics by Katrina Luckenbach and Matthew Vieira as
undergraduate research students~\cite[p.~13]{Hopkins}. Background on
Rado numbers is in Landman and Robertson's book~\cite[Chapter~9]{LR}.

\subsection{What is new}
Hopkins's proof of the case $\ell=1$ and the descent below both start
from the solution $(1,1,a+a\ell)$. For $\ell=1$ a single descent
already reaches a contradiction. For $\ell\ge2$ the descent stops at a
residue $r\ge2$, and the new step is the second stage: two families of
solutions that make color~$1$ invariant under shifts by $\ell$ on the
interval $[1,a\ell(r+1)]$. The lower-bound coloring is the same for all
$\ell$, and it is the hypothesis $a\ge2$ that makes it a golden
coloring.

\section{Proof overview}
Write $s=a(\ell+1)$, so the claimed value is $s^2$, and let
$b=a\ell$. Every positive solution of $ax+by=z$ has $z=a(x+\ell y)$,
a multiple of $a$ that is at least $s$. Coloring exactly these numbers
with~$1$ forces any monochromatic solution to have $x,y\ge s$, hence
$z\ge s^2$; this is the lower bound (Section~\ref{sec:lower}).

For the upper bound (Section~\ref{sec:upper}) take a golden coloring of
$[1,s^2]$ without monochromatic solutions. Because $1$ has color~$0$, the
solution $(1,1,s)$ gives $s$ color~$1$. The identity
$a(t-\ell)+b(t+1)=st=at+bt$ moves color~$1$ from $t$ down to $t-\ell$,
so color~$1$ descends from $s$ to its least positive residue $r$
modulo $\ell$. Then $r\ge2$, and both neighbors of $r$ have color~$0$.
Solutions with $y\in\{r-1,r,r+1\}$ propagate color~$1$ by $\pm\ell$
on $[1,T]$, where $T=b(r+1)$. Whatever the color of $\ell$, propagation
contradicts one of the solutions $(\ell,1,2b)$ and $(\ell,r,T)$. All
entries stay inside $[1,s^2]$ because $r\le\min(a,\ell)$.

\section{The lower bound}\label{sec:lower}
\begin{lemma}\label{lem:lower}
Let $a\ge2$, $\ell\ge1$ and $s=a(\ell+1)$. Color $t$ with $1$ exactly
when $a$ divides $t$ and $t\ge s$. For every $M$ with $1\le M<s^2$, the
restriction to $[1,M]$ is a golden coloring with no monochromatic
solution of $ax+a\ell y=z$.
\end{lemma}
\begin{proof}
Since $a\ge2$, the integer $1$ is not a multiple of $a$, and two
consecutive integers are never both multiples of $a$; so the coloring
is golden. Let $(x,y,z)$ be a solution with entries in $[1,M]$. Then
$z=a(x+\ell y)$ is a multiple of $a$ and $z\ge a(1+\ell)=s$, so $z$ has
color~$1$. If $x$ and $y$ also had color~$1$, then $x,y\ge s$ and
$z\ge a(s+\ell s)=s^2>M$, which is impossible.
\end{proof}

For $a=1$ every integer $t\ge s$ receives color~$1$, so the coloring is
not golden; this is where the hypothesis $a\ge2$ enters.

\section{The upper bound}\label{sec:upper}
Fix $a,\ell\ge1$, put $s=a(\ell+1)$ and $b=a\ell$, and let $\Delta$ be a
golden coloring of $[1,s^2]$ with no monochromatic solution of
$ax+by=z$. We derive a contradiction, which proves that every golden
coloring of $[1,s^2]$ contains a monochromatic solution. Two facts are
used throughout. If $(x,y,z)$ is a solution in $[1,s^2]$ and
$\Delta(x)=\Delta(y)$, then $\Delta(z)\ne\Delta(x)$; and if
$\Delta(y)=\Delta(z)$, then $\Delta(x)\ne\Delta(y)$. Also, a number of
color~$1$ has both neighbors in $[1,s^2]$ of color~$0$.

Let $r$ be the integer with $1\le r\le\ell$ and $r\equiv a\pmod\ell$,
that is, $r=((a-1)\bmod \ell)+1$. Then $r\le a$, so
$r\le\min(a,\ell)$, and $s=a+b\equiv a\equiv r\pmod\ell$.

\begin{lemma}[Descent]\label{lem:descent}
$\Delta(s)=1$, and $\Delta(t)=1$ with $\ell<t\le s$ implies
$\Delta(t-\ell)=1$. Consequently $\Delta(r)=1$, $r\ge2$ (so
$\ell\ge2$), and $\Delta(r-1)=\Delta(r+1)=0$.
\end{lemma}
\begin{proof}
The solution $(1,1,s)$ has $\Delta(1)=0$, so $\Delta(s)=1$. Let
$\ell<t\le s$ with $\Delta(t)=1$. The solution $(t,t,st)$, with
$st\le s^2$, gives $\Delta(st)=0$, and $\Delta(t+1)=0$ because the
coloring is golden. Since
\[
a(t-\ell)+b(t+1)=at+bt=st,
\]
the solution $(t-\ell,t+1,st)$ has $\Delta(t+1)=\Delta(st)=0$, so
$\Delta(t-\ell)=1$. Starting from $s$ and stepping down by $\ell$ while
the current value exceeds $\ell$, we stay in the residue class of $s$
modulo $\ell$ and stop at $r$; hence $\Delta(r)=1$. As $\Delta(1)=0$, we
get $r\ge2$, and the golden condition gives $\Delta(r\pm1)=0$.
\end{proof}

In particular, for $\ell=1$ we have $r=1$, and Lemma~\ref{lem:descent}
is already a contradiction. From now on $\ell\ge2$. Put
\[
T=b(r+1)=a\ell(r+1),\qquad K=a(\ell+1)^2,
\]
so that $aK=s^2$. Since $r\le\min(a,\ell)$,
\begin{equation}\label{eq:range}
T+\ell r=a\ell r+a\ell+\ell r\le a\ell^2+2a\ell<K,\qquad 3b\le T.
\end{equation}

\begin{lemma}[Shifts]\label{lem:shift}
Let $1\le t$ and $t+\ell\le T$. Then $\Delta(t)=1$ if and only if
$\Delta(t+\ell)=1$.
\end{lemma}
\begin{proof}
Let $u\le T$ be a positive integer with $\Delta(u)=1$. The solution
$(u,r,a(u+\ell r))$ has both inputs of color~$1$ by
Lemma~\ref{lem:descent}, and its output is at most $a(T+\ell r)<aK=s^2$
by~\eqref{eq:range}; so $\Delta(a(u+\ell r))=0$.

If $\Delta(t)=1$, apply this with $u=t$. The identity
$a(t+\ell)+b(r-1)=a(t+\ell r)$ gives the solution
$(t+\ell,r-1,a(t+\ell r))$, whose last two entries have color~$0$ (the
middle one by Lemma~\ref{lem:descent}, and $r-1\ge1$). Hence
$\Delta(t+\ell)=1$.

If $\Delta(t+\ell)=1$, apply it with $u=t+\ell$: the output
$a(t+\ell+\ell r)=at+b(r+1)$ has color~$0$. The solution
$(t,r+1,at+b(r+1))$ has its last two entries of color~$0$, so
$\Delta(t)=1$.
\end{proof}

\begin{proof}[Proof of Theorem~\ref{thm:main}]
By Lemma~\ref{lem:lower}, no positive $M<s^2$ has the forcing
property, so $\Rg(a,\ell a)\ge s^2$. For the upper bound, suppose
$\Delta$ is as above, with $\ell\ge2$. The multiples
$\ell,2\ell,\dots,T$ of $\ell$ all lie in $[1,T]$, so by
Lemma~\ref{lem:shift} they all have the same color.

If $\Delta(\ell)=0$, the solution $(\ell,1,2b)$ has
$\Delta(\ell)=\Delta(1)=0$, so the multiple $2b=2a\ell\le T$ of $\ell$
has color~$1$, a contradiction. If $\Delta(\ell)=1$, then $\Delta(T)=1$;
but $a\ell+br=b(r+1)=T$, so $(\ell,r,T)$ is a solution with
$\Delta(\ell)=\Delta(r)=1$ and therefore $\Delta(T)=0$, again a
contradiction. Thus every golden coloring of $[1,s^2]$ contains a
monochromatic solution, and $\Rg(a,\ell a)=s^2=a^2(\ell+1)^2$.
\end{proof}

\begin{remark}
The upper-bound argument does not use $a\ge2$: for every $a,\ell\ge1$,
each golden coloring of $[1,a^2(\ell+1)^2]$ contains a monochromatic
solution of $ax+a\ell y=z$. For $a=1$ and $\ell\ge2$ this bound is not
sharp, by Hopkins's Theorem~1.
\end{remark}

\section{Formal verification}
The theorem is proved in Lean~4 \cite{Lean} with Mathlib
\cite{Mathlib}, from the following definitions (the file
\texttt{Hopkins/Definitions.lean}). A coloring is a function
$c:\mathbb N\to\{\mathrm{false},\mathrm{true}\}$, with
$\mathrm{true}$ playing the role of color~$1$; only its values on
$[1,N]$ matter.
\begin{Verbatim}
def Valid (N : ℕ) (c : ℕ → Bool) : Prop :=
  c 1 = false ∧ ∀ i : ℕ, 1 ≤ i → i + 1 ≤ N →
    ¬ (c i = true ∧ c (i + 1) = true)

def Mono (a b N : ℕ) (c : ℕ → Bool) : Prop :=
  ∃ x y z : ℕ, 1 ≤ x ∧ x ≤ N ∧ 1 ≤ y ∧ y ≤ N ∧
    1 ≤ z ∧ z ≤ N ∧ a * x + b * y = z ∧ c x = c y ∧ c y = c z

def Forces (a b N : ℕ) : Prop :=
  ∀ c : ℕ → Bool, Valid N c → Mono a b N c

def ForcingSet (a b : ℕ) : Set ℕ := {N | 1 ≤ N ∧ Forces a b N}

noncomputable def R0phi (a b : ℕ) : ℕ := sInf (ForcingSet a b)
\end{Verbatim}
Thus \texttt{Valid N c} says that $c$ restricts to a golden coloring of
$[1,N]$, \texttt{Mono} asks for a monochromatic solution with all
entries in $[1,N]$ (with $x=y$ allowed), and \texttt{R0phi a b} is the
least positive $N$ with the forcing property, which is Hopkins's
$\Rg(a,b)$. The main theorems are
\begin{Verbatim}
theorem Hopkins.formula_isLeast (a ell : ℕ) (ha : 2 ≤ a) (hell : 1 ≤ ell) :
    IsLeast (ForcingSet a (a * ell)) (a ^ 2 * (ell + 1) ^ 2)
theorem Hopkins.formula (a ell : ℕ) (ha : 2 ≤ a) (hell : 1 ≤ ell) :
    R0phi a (a * ell) = a ^ 2 * (ell + 1) ^ 2
\end{Verbatim}
The first says that $a^2(\ell+1)^2$ has the forcing property and that no
smaller positive integer does, which is Theorem~\ref{thm:main}; the
second deduces the value of the infimum from it, so no convention about
the infimum of an empty set is involved. The development also proves
the two bounds separately: \texttt{Hopkins.upper} (forcing at
$a^2(\ell+1)^2$ for all $a,\ell\ge1$) and \texttt{Hopkins.lower}
(Lemma~\ref{lem:lower}, with the coloring \texttt{lowerColor}). As
checks on the definitions, it proves Hopkins's values $\Rg(1,1)=4$ and
$\Rg(1,2)=4$ from Hopkins's Theorem~1, the second showing the failure of the
printed formula at $a=1$, $\ell=2$.

The proof uses only Lean's standard logical axioms, with no added
axioms and no compiled evaluation. The README in the \texttt{lean} folder
accompanying this paper gives the versions, the commands that reproduce
the check and its axiom report.

\begin{thebibliography}{9}
\bibitem{GuoSun}
S.~Guo and Z.-W.~Sun, Determination of the two-color Rado number for
$a_1x_1+\dots+a_mx_m=x_0$, \emph{J. Combin. Theory Ser. A} \textbf{115}
(2008), 345--353. \url{https://doi.org/10.1016/j.jcta.2007.06.001}

\bibitem{Hopkins}
B.~Hopkins, Introducing shift-constrained Rado numbers,
\emph{Integers} \textbf{21A} (2021), \#A16, published August 23, 2021.
\url{https://math.colgate.edu/~integers/graham16/graham16.pdf}.
Reprinted in \emph{Number Theory and Combinatorics: A Collection in
Honor of the Mathematics of Ronald Graham}, De Gruyter, 2022.
\url{https://doi.org/10.1515/9783110754216-016}

\bibitem{HopkinsSchaal}
B.~Hopkins and D.~Schaal, On Rado numbers for
$\sum_{i=1}^{m-1}a_ix_i=x_m$, \emph{Adv. in Appl. Math.} \textbf{35}
(2005), 433--441. \url{https://doi.org/10.1016/j.aam.2005.05.003}

\bibitem{LR}
B.~M.~Landman and A.~Robertson, \emph{Ramsey Theory on the Integers},
2nd ed., American Mathematical Society, Providence, RI, 2014.

\bibitem{Lean}
L.~de Moura and S.~Ullrich, The Lean 4 theorem prover and programming
language, in \emph{Automated Deduction -- CADE 28}, Lecture Notes in
Comput. Sci. \textbf{12699}, Springer, 2021, 625--635.

\bibitem{LindMarcus}
D.~Lind and B.~Marcus, \emph{An Introduction to Symbolic Dynamics and
Coding}, 2nd ed., Cambridge University Press, Cambridge, 2021.

\bibitem{Mathlib}
The mathlib Community, The Lean mathematical library, in
\emph{Proceedings of the 9th ACM SIGPLAN International Conference on
Certified Programs and Proofs (CPP 2020)}, ACM, 2020, 367--381.
\end{thebibliography}
\end{document}
